\chapter*{INTRODUÇÃO}
\addcontentsline{toc}{chapter}{INTRODUÇÃO} % Adiciona a seção sem numeração no sumário
Este artigo é um registro que relata o conteúdo aprendido no Curso "Do Overleaf ao Git" lecionado por Lucas Alexandre Zick que teve como objetivo o aprendizado da linguagem de diagramação chamada \LaTeX e das ferramentas Git e GitHub como ferramentas complementares ao \LaTeX.

O curso foi composto por aulas teóricas e práticas, onde foi abordado desde a história e os conceitos básicos do \LaTeX, como o uso de pacotes, comandos e ambientes, até a criação de formatações próprias e o versionamento de código utilizando o Git e o GitHub.

O objetivo deste artigo é apresentar uma síntese dos conteúdos abordados em cada aula do curso e demonstrar a aplicação prática real dos conceitos ministrados no curso através deste documento. Este texto também foi pensado para servir como uma documentação para que futuros alunos possam conhecer e aprendam mais sobre o \LaTeX, pois pode ser útil para a trajetória acadêmica visto que sua grande maioria se mostra mais eficiente na escrita de documentos que necessitam de formatação complexa, como artigos científicos, e documentos mais longos que exigem uma organização mais robusta, como trabalhos acadêmicos e TCCs, quando comparado às outras ferramentas de escrita de documentos como o Word, e sendo principalmente mais eficiente quando necessitamos de referenciamento bibliográfico no documento.